\chapter{Data \& AI as Strategic Capability}\label{ch:data-and-ai-strategy}

% Scope: data as a strategic asset (data flywheel/network effects), AI build-vs-buy-vs-API,
%   AI governance & model risk, ESG as a strategic data problem (double materiality).
%   NEW chapter (closes the data-&-AI strategy gap + part of ESG depth).
% Operationalises VRIN (\cref{ch:advantage-and-disruption}) and moats (\cref{ch:growth}) for the AI context.
% Draft per house style: flowing prose + example/admonition set-offs; running-SaaS worked examples.
% This chapter is QUALITATIVE: no new eq: labels are defined here. All equations are \cref'd from existing labels.

\section{Data as a Strategic Asset}\label{sec:data-asset}
% complexity: 3

What makes a dataset competitively valuable? The intuition that ``more data is always better''
overstates the case: most data collections plateau in predictive accuracy long before they are
large enough to constitute a moat. The rarer phenomenon — a data flywheel — is a self-reinforcing
feedback loop in which a product attracts users, users generate data, data trains better predictions,
better predictions improve the product, and the improved product attracts more users. Each cycle
widens the gap against a competitor who must start the loop from scratch. The network-value logic
of \cref{eq:metcalfe} captures why this compounds: as the user base grows linearly, the combinatorial
space of behavioural signals grows quadratically, producing algorithmic improvements no static
dataset can replicate.

The strategic test for whether a data asset actually sustains advantage is the VRIN framework of
\cref{sec:rbv-vrin}: is it \emph{Valuable} (does it improve predictions customers care about?),
\emph{Rare} (can competitors acquire an equivalent corpus?), \emph{Inimitable} (is the data
path-dependent or historically contingent?), and \emph{Non-substitutable} (can a different
resource deliver the same product outcome?). For most internal operational data — server logs,
help-desk tickets, generic CRM activity — the answer to at least three of those four tests is no,
and the data advantage is illusory. A genuine data moat requires all four to be yes
\parencite{agrawal2018prediction}.

\begin{example}[SaaS multi-tenant benchmark data and the VRIN test]
A SaaS platform with \money{4000000} ARR collects anonymised usage telemetry across every tenant: feature
adoption rates, workflow completion times, error frequencies, and session-depth patterns. The product
team proposes a \emph{benchmarking} feature that shows each customer how their usage compares to
the peer cohort in their industry vertical. Run the VRIN test explicitly.

\textbf{Valuable?} Yes: benchmarks reduce customers' uncertainty about whether they are extracting
full value from the product, which increases stickiness. The improved retention lifts CLV directly
via \cref{eq:roic}: if monthly churn falls from \qty{0.65}{\percent} to \qty{0.55}{\percent} and
the discount rate is \qty{11}{\percent}, the present value of the customer base rises materially.

\textbf{Rare?} Yes: multi-tenant usage data at this granularity requires a live, consented customer
base at scale. A new entrant cannot synthesise it; they must earn it year by year.

\textbf{Inimitable?} Yes: the corpus embeds historical context — seasonal patterns, version-migration
behaviour, industry-specific onboarding paths — that took years to accumulate. Replicating the
algorithm is straightforward; replicating the training set is not.

\textbf{Non-substitutable?} Yes: no publicly available dataset delivers per-account, real-time
SaaS workflow telemetry at the feature level. Survey benchmarks and analyst reports are too coarse
and too lagged to serve the same purpose.

All four tests pass. This data asset sustains ROIC\,>\,WACC (\cref{eq:roic}) beyond any single
feature's life, because each new cohort of customers enriches the corpus and strengthens the
benchmark signal that differentiates the product.
\end{example}

\begin{admonition}{Warning}
Data advantage degrades under data-portability regulation. Article~20 of the GDPR grants individuals
the right to receive their personal data in a portable format and to transmit it to a competing
provider. The EU Digital Markets Act extends a structurally similar obligation to gatekeepers. A
data moat whose feedstock is personal data rather than aggregated behavioural signal can be
eroded by portability rights faster than a new product cycle can refresh the competitive advantage.
Audit whether your data flywheel survives a portability regime before pricing it into your
strategic plan.
\end{admonition}

The data flywheel and the network-effects model of \cref{eq:metcalfe} share the same structural
logic: value accrues non-linearly as the user base grows, but only when the marginal datum
genuinely improves the core output. Both mechanisms connect back to the sustainable ROIC spread of
\cref{eq:economic-profit}: the data asset is valuable only insofar as it generates a product
improvement that customers will pay to retain. \textcite{agrawal2018prediction} develop the
economics of prediction as a falling-cost input and its implications for competitive advantage.

\section{AI Build vs.\ Buy vs.\ API}\label{sec:ai-build-buy}
% complexity: 4

Every AI investment decision reduces to a version of the same question: at what level of the
stack should the firm own the capability rather than purchase it? The economics of AI have shifted
dramatically since the early 2020s: the cost of running inference on a powerful general-purpose
model has fallen by orders of magnitude, making the API option the correct default for most tasks
and raising the bar for when building proprietary capability is justified. The decision is
best framed as a two-dimensional mapping: how proprietary is the firm's training data, and how
task-specific is the application?

The 2×2 in \cref{fig:ai-build-buy} organises the four quadrants. When both data and task are
generic, buying an API wrapper is correct: the economics of commodity inference dominate and no
sustainable moat exists. When the task is generic but the firm owns proprietary data, fine-tuning
a base model on that data can provide mild differentiation at moderate cost. When neither data
nor task is proprietary, orchestrating a vendor SaaS layer is usually the right answer: the firm
captures the productivity gain without the maintenance burden. Only in the upper-right quadrant —
high proprietary data, high task specificity — does a full build create a defensible asset; this
is the regime of the VRIN-positive data flywheel in \cref{sec:data-asset}.

\begin{figure}[htbp]
  \centering
  \includegraphics[width=0.62\linewidth]{ai-build-buy}
  \caption{AI procurement framework. Horizontal axis: task specificity (generic to proprietary
           use case). Vertical axis: proprietary data advantage (low to high). The build quadrant
           is the only one where full ownership generates a VRIN-positive asset;
           the other three favour buying or orchestrating external capability.}
  \label{fig:ai-build-buy}
\end{figure}

The financial discipline is \cref{eq:npv}: build is justified only when the present value of
the differentiated benefit exceeds the present value of the total cost of ownership — upfront
training, ongoing inference, and the maintenance cost that compounds as model drift requires
retraining. The buy option has a near-zero upfront cost and a fully known per-unit price; it
carries no moat but also no engineering debt.

\begin{example}[AI report summaries: API vs.\ fine-tune]
The SaaS firm is evaluating two options for an AI-generated narrative summary feature that
appears in each customer's monthly analytics report.

\textbf{Option A — API wrapper:} Use a commodity large-language model via API at
\money{0.002} per thousand tokens. At typical report volume, this yields approximately
\money{10} per month in inference costs (\money{120} per year). Implementation takes two
engineering weeks; no moat is created, because any competitor can call the same endpoint.

\textbf{Option B — Fine-tune on usage data:} Adapt a base model on the firm's proprietary
multi-tenant telemetry corpus. One-time training cost: \money{15000}. Ongoing inference
and retraining overhead: \money{2000} per month (\money{24000} per year). The fine-tuned
model generates summaries that cite the customer's own benchmarks and workflow patterns,
mild VRIN because the training data passes three of the four tests in \cref{sec:data-asset}.

The decision horizon is three years at the firm's \qty{11}{\percent} WACC. The incremental
cost of Option B over Option A is the \money{15000} upfront plus the annual differential
of \money{23880} (\money{24000} minus \money{120}), discounted at \qty{11}{\percent}:
\[
  \Delta\mathrm{Cost} \;=\; \money{15000}
    \;+\; \money{23880}\times
    \underbrace{\sum_{t=1}^{3}\frac{1}{(1.11)^{t}}}_{=\;2.4437}
  \;\approx\; \money{73360}.
\]

The benefit of the fine-tuned model flows through retention: a subset of approximately 600
report-active customers responds to the richer summaries, and the differentiation reduces
churn for that cohort. The NPV is positive only if the CLV lift from this retention
exceeds roughly \money{50} per customer, because
\[
  \underbrace{600\;\times\;\money{50}}_{\text{customers}\;\times\;\Delta\mathrm{CLV}}
  \;\times\; 2.4437
  \;\approx\; \money{73300},
\]
which barely clears the incremental cost. If the feature produces only modest
stickiness — say \money{30}/customer — the NPV is firmly negative and Option A is correct.
The implication: do not build unless a plausible mechanism exists by which the fine-tuned
model creates \emph{meaningful} differentiated retention, not just a marginal quality
improvement.
\end{example}

\begin{admonition}{Note}
The total cost of ownership for a custom AI system is systematically underestimated at the
prototype stage. The ``last five percent'' — security hardening, latency SLAs, drift monitoring,
retraining pipelines, and compliance logging — routinely represents the majority of lifetime
engineering spend. An \cref{eq:npv} analysis that captures only upfront training cost will
always favour building; one that amortises the full maintenance burden typically does not.
Model the full cost before committing.
\end{admonition}

The build-versus-API choice is structurally a moat test: only build when the output of building
clears the VRIN hurdle of \cref{sec:rbv-vrin} and produces a sustained ROIC\,>\,WACC spread.
In the SaaS context, the proprietary-data axis of \cref{fig:ai-build-buy} maps directly onto
the data-flywheel logic of \cref{sec:data-asset}: without proprietary data, fine-tuning buys
a temporary quality improvement, not a moat. \textcite{agrawal2018prediction} analyse the
complementarity between cheap prediction and scarce human judgment as the governing logic of
AI investment prioritisation.

\section{AI Governance \& Model Risk}\label{sec:ai-governance}
% complexity: 3

Deploying AI in a commercial context introduces a category of risk that does not appear in a
traditional software risk register. A conventional software bug produces a wrong output in a
known and testable way; a generative AI model produces a wrong output \emph{confidently}, in
natural language, in a way that looks correct and may not be caught before it affects a decision.
The risk taxonomy has four principal dimensions: \emph{bias} (systematic error correlated with
protected or sensitive attributes), \emph{hallucination} (confident generation of factually
incorrect content), \emph{intellectual property} (training-data provenance and output ownership),
and \emph{regulatory} (jurisdiction-specific prohibitions on high-risk AI use cases, exemplified
by the EU AI Act's risk classification system).

The governance question is an agency problem in the sense of \cref{sec:agency-theory}: the model
is an agent whose objective function (next-token probability) is not perfectly aligned with the
firm's objective (accurate, useful output). The principal — the firm deploying the model — must
design oversight mechanisms to catch misalignment before it causes harm. This is precisely the
control-system logic that \cref{sec:agency-theory} applies to human agents: incentive design,
monitoring, and bounded discretion. For AI systems, bounded discretion becomes \emph{human-in-the-loop}
for high-stakes outputs, and monitoring becomes automated output validation.

The expected-value framework of \cref{sec:expected-value} provides the correct risk-triage tool.
For each AI-assisted output category, the firm should estimate the probability of a material error,
the value at risk if that error propagates undetected, and the cost of the human-review layer.
Governance resources should be concentrated where the product of probability and value-at-risk
is largest.

\begin{example}[Revenue projection hallucination and expected value at risk]
The SaaS firm has deployed the AI summary feature from the previous section. In one case, the
model emits a confidently-worded revenue projection in a customer's report that overstates
the customer's projected growth by \qty{40}{\percent}. The customer uses the projection to
plan headcount, discovers the error three months later, and churns.

The expected value at risk of a single such incident is the CLV of a churned customer. At the
firm's standard parameters (\cref{eq:clv}: \money{252}/yr contribution margin, \qty{11}{\percent}
WACC, \qty{3}{\percent} growth), CLV\,=\,\money{3150}. If the model's rate of material
projection errors is estimated at one per 1{,}050 report runs per month, and each error
has a \qty{50}{\percent} probability of triggering churn, then the expected monthly churn
cost attributable to the feature is \(\money{3150}\times 0.50\approx\money{1575}\).

A governance design that catches the error before customer delivery requires one analyst to
review all financial-output summaries for approximately 30 minutes per week — a cost well
below \money{1575}/month. More precisely: the \emph{worst-case} scenario of three simultaneous
uncaught errors producing three churned customers creates an expected value at risk of
\(3\times\money{3150}=\money{9450}\), which is the financial floor for the governance
investment. Human-in-the-loop review of financial outputs is clearly positive-NPV. For
operational summaries (feature adoption, workflow completions) where an error is visible to
the customer but carries low financial consequence, automated guardrails and downstream
correction are sufficient.
\end{example}

\begin{admonition}{Warning}
``AI-generated'' is not a legal defence. If a model emits a materially misleading financial
projection, an incorrect medical recommendation, or a defamatory statement in a customer-facing
document, the liability attaches to the firm that deployed the model, not to the model or its
vendor. Contractual disclaimers shift some commercial risk but do not displace statutory or
tortious liability. This is the regulatory dimension of the VRIN test: a non-substitutable data
moat may simultaneously be a non-substitutable regulatory exposure.
\end{admonition}

The governance framework connects to attribution in \cref{ch:attribution-and-measurement}: just
as attribution models require a closed feedback loop between model output and ground-truth
outcome, AI governance requires a systematic mechanism for flagging model errors, routing them
to human review, and using the resulting signal to retrain or constrain the model.
\textcite{jensen1976agency} establish the principal-agent structure underlying these control
problems; the AI context is a new instantiation of the alignment challenge they analyse.

\section{ESG as a Strategic Data Problem}\label{sec:esg-data}
% complexity: 3

Environmental, social, and governance reporting is usually framed as a compliance obligation.
Its strategic significance is different and larger: ESG is a \emph{measurement problem} with
the same structure as attribution (\cref{ch:attribution-and-measurement}). Both require the
firm to trace multi-period causal effects through indirect channels, aggregate data from
heterogeneous sources, and report on metrics that affect capital allocation decisions. Both
are subject to misattribution, cherry-picking, and model uncertainty. The discipline required
to measure carbon footprint accurately is exactly the discipline required to measure marketing
incrementality accurately: clear counterfactuals, boundary conditions, and auditability.

The governing framework is \emph{double materiality}, codified in ISSB standards S1 and S2
and in the European Sustainability Reporting Standards (ESRS). Double materiality requires
a firm to disclose sustainability topics from two perspectives simultaneously. \emph{Financial
materiality} (outside-in) asks how ESG factors — physical climate risk, supply-chain
disruptions, talent shortages, regulatory tightening — affect the firm's enterprise value
and cash flows. \emph{Impact materiality} (inside-out) asks how the firm's operations affect
people and the environment, irrespective of whether those impacts currently flow back to the
income statement. A topic is reportable if it clears the threshold on \emph{either}
dimension; the union, not the intersection, determines the disclosure boundary.

The stakeholder-value logic of \cref{sec:stakeholder-value} provides the strategic rationale.
\textcite{freeman2010strategic} argues that sustainable value creation requires managing
obligations to employees, customers, suppliers, communities, and investors simultaneously.
Double materiality is the reporting architecture that makes that obligation measurable and
auditable.

\begin{example}[Scope 1/2/3 and Series~A diligence for the SaaS firm]
The SaaS firm with \money{4000000} ARR operates on cloud infrastructure. Scope~1 and
Scope~2 emissions (direct combustion and purchased electricity) are near-zero: the firm
owns no physical plant, and its cloud vendor reports renewable-energy coverage exceeding
\qty{90}{\percent}. On a pure financial-materiality analysis, carbon risk is not material
at this stage: no carbon tax or physical climate event is plausibly large enough to affect
a \money{4000000} ARR business within a five-year horizon.

The material ESG risks are different in character. Scope~3 emissions in the upstream supply
chain (hardware manufacturing, data-centre embodied carbon) are difficult to measure but
are increasingly expected in investor diligence. On the social dimension, talent equity —
pay equity across gender and ethnicity, stock-option distribution — is a genuine strategic
risk: at \money{120000}/month in acquisition spend and a 200 customer/month target, the
firm depends on a small, high-skill team; equity disputes or attrition create direct
operating risk. Neither item is financially material against \money{4000000} ARR in
isolation, but \emph{impact materiality} requires disclosure for Series~A diligence under
ISSB-aligned frameworks: institutional and strategic investors applying ISSB S1/S2 will
expect a double-materiality assessment, even from a company well below the CSRD reporting
threshold.

The economic-profit identity of \cref{eq:economic-profit} is the correct numerator for
evaluating ESG spend: any ESG investment that does not produce a reduction in WACC
(through lower regulatory risk or lower cost of capital) or an increase in NOPAT (through
improved talent retention or reduced churn) does not create economic profit and should be
evaluated accordingly. At \money{4000000} ARR, most ESG spend falls into the
cost-of-capital channel — the reduction in fundraising friction — rather than the NOPAT
channel.
\end{example}

\begin{admonition}{Note}
A high ESG score does not imply ROIC\,>\,WACC. The relationship between ESG performance
and economic profit runs through two specific mechanisms: reduced cost of capital (lower
risk premium from investors screening ESG exposure) and reduced operating risk (lower
probability of regulatory penalties or talent disruption). Neither mechanism is automatic;
both are measurable. The warning against conflating stakeholder optics with shareholder
value is developed fully in \cref{ch:value-creation} and is not restated here.
\end{admonition}

The attribution parallel is more than structural: the methodological tools are largely the
same. Measuring Scope~3 emissions requires a value-chain model with counterfactual baselines —
what would the emissions have been absent the firm's operations? — which is the same question
attribution models ask about media spend. Building ESG measurement capability therefore
compounds with attribution measurement capability; the two investments share infrastructure.
\textcite{freeman2010strategic} and \textcite{jensen1976agency} together define the
governance and value-creation tension that double materiality is designed to resolve.
